%==================================================================
%  main.tex  --  IEEE S&P 2027 submission
%  IEEE Computer Society conference template (IEEEtran.cls V1.8b,
%  [conference,compsoc]), anonymized for double-blind review.
%  The previous LNCS version is kept in main_llncs_backup.tex.
%==================================================================
\documentclass[conference,compsoc]{IEEEtran}

%------------------------------------------------------------------
%  Packages
%------------------------------------------------------------------
\usepackage[T1]{fontenc}
\usepackage[utf8]{inputenc}
\usepackage{graphicx}
\usepackage{amsmath,amssymb}  % amsthm은 로드하지 말 것 (IEEEtran의 IEEEproof 사용)
\interdisplaylinepenalty=2500 % IEEEtran 권장: amsmath 사용 시 수식 내 페이지 나눔 허용
\usepackage{booktabs}
% IEEE Computer Society는 [1]-[4] 대신 [1], [2], [3], [4] 형식 (nocompress)
\ifCLASSOPTIONcompsoc
  \usepackage[nocompress]{cite}
\else
  \usepackage{cite}
\fi
\usepackage{algorithm}
\usepackage{algpseudocode}
\usepackage{xcolor}           % color 대신 xcolor (TikZ와 공용)
\usepackage{tikz}             % Fig. 1 (Technical Overview)
\usetikzlibrary{arrows.meta,shadows}
\usepackage{hyperref}         % 가급적 마지막에 로드
\hypersetup{pdfauthor={},pdfcreator={},pdfproducer={}} % 익명 심사: PDF 메타데이터 비움

% IEEEtran은 theorem 환경을 제공하지 않으므로 직접 정의
\newtheorem{theorem}{Theorem}

%------------------------------------------------------------------
%  Notation / figure colors
%------------------------------------------------------------------
\newcommand{\enc}[1]{[\![#1]\!]}          % ciphertext of #1
% Section 5 표기 (Table 5가 Section 4 파일에서 정의되므로 preamble에 둠)
\providecommand{\MinPts}{\mathit{MinPts}}
\providecommand{\kmax}{k_{\max}}
\providecommand{\epsn}{\varepsilon_n}
\providecommand{\pk}{\mathsf{pk}}
\providecommand{\sk}{\mathsf{sk}}
\providecommand{\evk}{\mathsf{evk}}
\providecommand{\ct}{\mathsf{ct}}
\definecolor{figNavy}{HTML}{0B1E31}       % 원본 그림의 박스 색
\definecolor{figGold}{HTML}{C8871A}       % 암호문 전송 화살표·자물쇠

%------------------------------------------------------------------
%  Folder names  (Overleaf 폴더 이름과 정확히 일치해야 함)
%------------------------------------------------------------------
\newcommand{\dirI}{Section1}     % Introduction (+ Related Work, Technical Overview as subsections)
\newcommand{\dirII}{Section2}    % Preliminaries
\newcommand{\dirIII}{Section3}   % FHE-DBSCAN
\newcommand{\dirIV}{Section4}    % Implementation results
\newcommand{\dirV}{Section5}     % Security analysis, Conclusion

% 파일이 없으면 컴파일을 멈추지 않고 PDF에 "Missing file" 상자를 표시
\newcommand{\inputsection}[1]{%
  \IfFileExists{#1.tex}{\input{#1}}{%
    \par\noindent\fbox{\ttfamily Missing file: \detokenize\expandafter{#1}.tex}\par
    \typeout{*** Missing file: #1.tex ***}}}

\graphicspath{{\dirI/Figures/}}

\hyphenation{Neigh-bor}

\begin{document}

%------------------------------------------------------------------
%  Front matter  (anonymous submission: no names or affiliations)
%------------------------------------------------------------------
\title{Non-Interactive Privacy-Preserving DBSCAN\\ with Homomorphic Encryption}
\author{\IEEEauthorblockN{Anonymous Author(s)}}
\maketitle

% IEEE: 초록에는 수식, 특수기호, 인용을 넣지 않음
\begin{abstract}
DBSCAN clusters data by density: it finds clusters of arbitrary shape,
labels noise explicitly, and needs no prior number of clusters. Evaluating
it on encrypted data is difficult because the algorithm is built on
comparisons, which homomorphic encryption does not support natively.
Existing privacy-preserving DBSCAN schemes therefore return comparisons to
the data owner or split the computation among non-colluding servers. We
present, to our knowledge, the first non-interactive DBSCAN over
homomorphic encryption. The client uploads its data once, encrypted under
the Cheon--Kim--Kim--Song (CKKS) scheme, and receives a single ciphertext of cluster labels. In between, a single
server evaluates the entire algorithm on ciphertexts. Instead of comparing
every pair of points, the server compares each point only within a range
that contains all of its neighbors, which greatly reduces the computation.
Since core-point detection and label propagation operate on discrete
values, we design their comparisons so that every input that decides the
output lies outside the interval where the polynomial approximation of the
sign function gives no accuracy guarantee. On six benchmark datasets and
two real datasets from finance and healthcare, the
encrypted result is equivalent to plaintext DBSCAN under the same
border-point rule, while the server learns nothing beyond a few public
parameters.
\end{abstract}

\IEEEpeerreviewmaketitle

%------------------------------------------------------------------
%  Folder 1: Introduction / Related Work / Technical Overview
%------------------------------------------------------------------
\inputsection{\dirI/Introduction}
\inputsection{\dirI/RelatedWork}
\inputsection{\dirI/TechnicalOverview}

%------------------------------------------------------------------
%  Folder 2: Preliminaries            (\label{sec:prelim})
%------------------------------------------------------------------
\inputsection{\dirII/Preliminaries}

%------------------------------------------------------------------
%  Folder 3: Proposed Method          (\label{sec:method})
%------------------------------------------------------------------
\inputsection{\dirIII/ProposedMethod}        % 절 제목 + 도입
\inputsection{\dirIII/ClientPreprocessing}   % 3.1 (RoundCount 포함)
\inputsection{\dirIII/NeighborCheck}         % 3.2 머리말 + 3.2.1
\inputsection{\dirIII/CoreCheck}             % 3.2.2
\inputsection{\dirIII/LabelPropagation}      % 3.2.3
\inputsection{\dirIII/Correctness}           % 3.2.4
\inputsection{\dirIII/PostProcessing}        % 3.3
\inputsection{\dirIII/Figures/FullProtocol}  % Algorithm 4 (전체 프로토콜)

%------------------------------------------------------------------
%  Folder 4: Implementation results   (\label{sec:eval})
%------------------------------------------------------------------
\inputsection{\dirIV/Implementation}

%------------------------------------------------------------------
%  Folder 5: Security analysis / Conclusion
%------------------------------------------------------------------
\inputsection{\dirV/SecurityAnalysis}
\inputsection{\dirV/Conclusion}

%------------------------------------------------------------------
%  Acknowledgments  (익명 심사: 연구비·동료 이름 등 신원 정보 금지)
%  IEEE compsoc 템플릿 형식에 따라 번호 없는 \section* 사용
%------------------------------------------------------------------
\section*{Acknowledgments}
Generative AI was used for editorial purposes in this manuscript, and all
outputs were inspected by the authors to ensure accuracy and originality.

%------------------------------------------------------------------
%  References
%------------------------------------------------------------------
\bibliographystyle{IEEEtran}
\bibliography{references}

%------------------------------------------------------------------
%  Appendices  (13쪽 이후: S&P 규정상 appendix로 명시)
%------------------------------------------------------------------
\appendices
\inputsection{Appendix/CorrectnessDetails}

\end{document}
